\section{高考考频分析}

\subsection{近五年考点分布（2021--2025 新课标Ⅰ卷）}

\begin{center}
\begin{tabular}{c|c|c|c|c|c}
\textbf{知识模块} & \textbf{2021} & \textbf{2022} & \textbf{2023} & \textbf{2024} & \textbf{2025} \\ \hline
集合 & 5 & 5 & 5 & 5 & 5 \\
复数 & 5 & 5 & 5 & 5 & 5 \\
不等式 & 0 & 5 & 0 & 0 & 5 \\
函数与导数 & 22+5 & 22+5 & 22+10 & 22+10 & 22+5 \\
三角函数与解三角形 & 10 & 10 & 12+5 & 10+5 & 10 \\
数列 & 12 & 12+5 & 10 & 10 & 12 \\
平面向量 & 5 & 5 & 0 & 5 & 5 \\
立体几何 & 12+5 & 12 & 12+5 & 12 & 12 \\
解析几何 & 12+5 & 12+10 & 12+5 & 12+5 & 12+10 \\
概率统计 & 12+5 & 12+5 & 12+5 & 12 & 12+5 \\
排列组合与二项式 & 5 & 5 & 5 & 5 & 5 \\ \hline
\textbf{总计} & 150 & 150 & 150 & 150 & 150
\end{tabular}
\end{center}

\textit{注：数字为分值，"22+5" 表示解答题 22 分 + 选填 5 分。}

\subsection{必考板块（每年必出）}

\[
\begin{array}{c|c|c}
\text{模块} & \text{年均分值} & \text{考查形式} \\ \hline
\text{函数与导数} & 25\text{--}30 分 & \text{选填 1--2 道 + 压轴解答题} \\
\text{解析几何} & 15\text{--}22 分 & \text{选填 1 道 + 解答 21 题} \\
\text{立体几何} & 12\text{--}17 分 & \text{选填 1 道 + 解答 19 题} \\
\text{概率统计} & 12\text{--}17 分 & \text{选填 1 道 + 解答 20 题} \\
\text{三角函数} & 10\text{--}17 分 & \text{选填 1 道 + 解答 17/18 题} \\
\text{数列} & 10\text{--}17 分 & \text{选填 1 道 + 解答 17/18 题} \\
\text{集合、复数、排列组合} & 5 分 & \text{仅选填}
\end{array}
\]

\subsection{题型分布特征}

\textbf{单选题分布}（8 题 × 5 分 = 40 分）：
\begin{itemize}
  \item 第 1--2 题：集合、复数（送分）
  \item 第 3--4 题：向量、排列组合、不等式等基础
  \item 第 5--6 题：三角、数列、函数性质中档题
  \item 第 7--8 题：函数综合、导数、解析几何压轴选填
\end{itemize}

\textbf{多选题分布}（4 题 × 5 分 = 20 分）：
\begin{itemize}
  \item 第 9 题：中档（常考函数性质、统计图表等）
  \item 第 10 题：中档偏上
  \item 第 11--12 题：难题（常考抽象函数、解析几何、立体几何动态问题）
\end{itemize}

\textbf{填空题分布}（4 题 × 5 分 = 20 分）：
\begin{itemize}
  \item 第 13--14 题：基础计算（二项式系数、数列、向量等）
  \item 第 15--16 题：中档偏难（解三角形实际应用、圆锥曲线、导数）
\end{itemize}

\textbf{解答题分布}（6 题共 70 分）：

\[
\begin{array}{c|c|c}
\text{题号} & \text{分值} & \text{常考内容} \\ \hline
17 & 10 & \text{数列或三角（基础，必拿）} \\
18 & 12 & \text{三角或数列（中档）} \\
19 & 12 & \text{立体几何（向量法建系，中档）} \\
20 & 12 & \text{概率与统计（情境题，中档）} \\
21 & 12 & \text{解析几何（椭圆/双曲线，偏难）} \\
22 & 12 & \text{函数与导数（压轴，第一问必拿）}
\end{array}
\]

\subsection{命题趋势分析}

\begin{enumerate}
  \item \textbf{反套路化}：近年不再固定题号对应固定知识模块，数列和三角交替出现在 17/18 题
  \item \textbf{情境题增多}：概率统计题常结合实际情境（投篮、质检、环保等），理解题意是关键
  \item \textbf{导数难度增加}：导数压轴题从纯技巧转向分析能力，隐零点、同构、极值点偏移是高频考点
  \item \textbf{解析几何回归本质}：淡化繁杂计算，强化学对几何关系的代数转化能力
  \item \textbf{新定义问题}：近年出现"新定义"题型，考察现场学习能力和知识迁移能力
  \item \textbf{多选题策略更重要}：多选题分值高、容错低，"宁少勿滥"的做题策略直接决定得分
\end{enumerate}

\subsection{复习优先级建议}

\[
\begin{array}{c|c|c}
\text{优先级} & \text{模块} & \text{理由} \\ \hline
\textbf{第一梯队} & \text{函数与导数、解析几何、立体几何} & \text{分值最高，拉开差距的关键} \\
\textbf{第二梯队} & \text{三角、数列、概率统计} & \text{必考，且是解答题得分基本盘} \\
\textbf{第三梯队} & \text{向量、排列组合、不等式} & \text{选填必考，但分值少，快速突破} \\
\textbf{第四梯队} & \text{集合、复数} & \text{送分题，不要丢分}
\end{array}
\]

\textbf{目标 90 分}：重点突破第一梯队基础解法 + 第二梯队全部 + 确保第四梯队不丢分。

\textbf{目标 120 分}：在第一梯队基础上，强化第一梯队的综合题 + 第三梯队秒杀技巧。

\textbf{目标 140+}：全模块无死角，重点攻关导数压轴、解析几何综合与多选题难题。
